\chapter{Method C --- Streaming mode}
\label{sec:stream}

\begin{center}
\small
\begin{tabular}{@{}ll@{}}
\toprule
Notebook cell & Step 4C (analysis: Step 5C) \\
Code path & \co{prog.stream(total_reps=N)} --- a generator \\
One packet returns & 38--500 reps, continuously, as they arrive \\
Measured rate & \textbf{1041.7\,Hz} (or 4167\,Hz with no host binning) \\
Rate limited by & the FPGA cadence alone \\
Duty cycle & \textbf{100.0\%} \\
Measured floor & 240\,\etaunit \\
Open defects & \Open{} \textbf{floor is $\sim$6$\times$ burst} (\S\ref{sec:streamnoise});
              25\% PL droop (\S\ref{sec:droop}) \\
\bottomrule
\end{tabular}
\end{center}

\section{Mechanics}

One \co{config_all}, one \co{start_readout(total_reps)}, then nothing but draining
the streamer queue with \co{poll_data()} in a loop. The board-side worker thread
transfers completed shots out of the accumulated buffer autonomously while the
tProc keeps running.

Because nothing is re-armed mid-run:

\begin{itemize}
  \item the replay race of Method B \textbf{cannot occur} --- there is no second
        configuration to race against;
  \item there is \textbf{no inter-burst gap} --- the FPGA never stops;
  \item timestamps follow the FPGA cadence \textbf{exactly}, rather than being
        reconstructed from host arrival times, which only record when the host
        happened to drain the queue.
\end{itemize}

\paragraph{\co{cfg.reps} is the run length.} The tProc loops \co{reps} times and
ends, so the rep count must be set before compilation and cannot change mid-run.
\co{stream()} raises if \co{cfg.reps != total_reps} rather than silently
truncating.

\paragraph{Host-side binning.} \co{STREAM_TARGET_HZ} sets an integer number of
reps to average on the host per output sample. Whole reps only, so the achieved
rate is the cadence divided by an integer rather than exactly the target: at
$\tau=\SI{120}{\micro\second}$, a 1000\,Hz target gives $\text{bin}=4$ and
therefore 1041.7\,Hz. Leftover reps carry into the next packet, so no rep is
dropped and no bin straddles a packet boundary.

\section{Timing budget, component by component}
\label{sec:streambudget}

\subsection*{Operating point: $\tau=\SI{120}{\micro\second}$, bin = 4 (2026-08-14)}

\begin{center}
\small
\begin{tabular}{@{}>{\raggedright\arraybackslash}p{4.3cm}rr>{\raggedright\arraybackslash}p{5.4cm}@{}}
\toprule
\hd{Component} & \hd{Per output sample} & \hd{Share} & \hd{How obtained} \\
\midrule
ADC integration ($4\times\SI{240}{\micro\second}$)
  & \textbf{\SI{960.0}{\micro\second}} & \textbf{100.0\%} & exact, Eq.~\ref{eq:recover}, $D$=147456 \\
Host RPC chain
  & \SI{0}{\micro\second} & 0\% & \textbf{paid once for the whole run} \\
Inter-packet gap
  & \SI{0}{\micro\second} & 0\% & the FPGA does not stop between packets \\
Timestamp jitter
  & \SI{0.0}{\micro\second} & 0\% & cadence-derived; min = max = \SI{960.0}{\micro\second} \\
\midrule
\textbf{TOTAL per output sample} & \textbf{\SI{960.0}{\micro\second}} & 100\% & measured spacing \\
Rate & \textbf{1041.67\,Hz} & & predicted 1041.67\,Hz --- \emph{exact} \\
Duty cycle & \textbf{100.0\%} & & 125\,000 reps $\times$ \SI{240}{\micro\second} in \SI{29.999}{\second} \\
\bottomrule
\end{tabular}
\end{center}

\begin{keybox}[This is what a clean timing budget looks like]
Every microsecond of the \SI{960}{\micro\second} sample period is ADC integration.
The measured sample spacing has \textbf{min = max = \SI{960.0}{\micro\second}} ---
zero jitter, because the timestamps come from the cadence rather than from arrival
times. Predicted and measured rate agree to the digit.

Compare Method A, where 79\% of the period is host overhead, and Method B, where
26\% is overhead and waste.
\end{keybox}

\section{Every parameter of the streaming cell, one at a time}
\label{sec:streamparams}

Step 4C has six knobs and three of them behave in ways that are not obvious from
their names. This section is the reference for all of them.

\subsection{\texorpdfstring{\co{STREAM_TARGET_HZ} and the bin}{STREAM\_TARGET\_HZ and the bin}}
\label{sec:bin}

You ask for a rate; the cell converts it into an integer number of reps to average
per output sample, and that integer --- not your request --- sets the rate:

\begin{equation}
\co{bin} = \max\!\left(1,\ \operatorname{round}\!\left(\frac{1}{\co{STREAM_TARGET_HZ}\times t_\text{rep}}\right)\right),
\qquad
f_\text{achieved} = \frac{1}{\co{bin}\times t_\text{rep}}.
\label{eq:bin}
\end{equation}

At $t_\text{rep}=\SI{240}{\micro\second}$, asking for \SI{1000}{\hertz} gives
$\co{bin}=\operatorname{round}(4.167)=4$ and therefore
\textbf{\SI{1041.7}{\hertz}}, not \SI{1000}{\hertz}. That is why every streamed
file in this project reports \SI{1041.7}{\hertz}: it is the nearest achievable
rate, and reps cannot be split.

Four things follow from the bin being a \emph{mean} of consecutive reps rather
than a decimation, and all four matter:

\begin{enumerate}
  \item \textbf{It averages, so it buys $\sqrt{\co{bin}}$ in $\sigma$.} Dropping
  3 of every 4 reps would not. At \co{bin} = 4 that is a factor of 2 --- the same
  free averaging Method A gets from filling the call floor (\S\ref{sec:freeavg}),
  obtained here without giving up any rate you were going to use.
  \item \textbf{It anti-aliases, at $1/(2\,\co{bin}\,t_\text{rep})$.} A boxcar of
  \co{bin} reps has its first null at $f_\text{achieved}$ and rolls off as
  $|\operatorname{sinc}|$ before it, so noise between $f_\text{achieved}/2$ and
  the raw \SI{4167}{\hertz} is suppressed rather than folded back in. This is the
  reason to bin in the acquisition cell rather than to record raw and decimate
  afterwards.
  \item \textbf{Bins never straddle a packet boundary.} Reps left over at the end
  of a packet are carried into the next one (\co{_pending} in the cell), so a bin
  is always \co{bin} consecutive reps of the same continuous run. Packet
  boundaries are an artefact of when the host drained the queue and mean nothing
  physical; a bin that straddled one would still be correct, but the carry makes
  the sample count exactly $\lfloor N_\text{reps}/\co{bin}\rfloor$ regardless of
  how the packets happened to land.
  \item \textbf{The timestamp is the bin centre.} \co{(first_rep + (i + 0.5) *
  bin) * cadence}. Using the bin's first rep would bias every sample early by
  $\co{bin}\,t_\text{rep}/2$ --- \SI{480}{\micro\second} at \co{bin} = 4.
\end{enumerate}

\begin{center}
\small
\begin{tabular}{@{}rrrl@{}}
\toprule
\hd{\co{bin}} & \hd{Per output sample} & \hd{Rate} & \hd{Note} \\
\midrule
1 & \SI{240}{\micro\second} & \textbf{4166.7\,Hz} & \co{STREAM_TARGET_HZ = None}; every rep kept \\
2 & \SI{480}{\micro\second} & 2083.3\,Hz & \\
\textbf{4} & \textbf{\SI{960}{\micro\second}} & \textbf{1041.7\,Hz} & what ran on 2026-08-14 and 2026-08-17 \\
8 & \SI{1920}{\micro\second} & 520.8\,Hz & \\
42 & \SI{10.08}{\milli\second} & 99.2\,Hz & equivalent averaging to Method A, at 100\% duty \\
\bottomrule
\end{tabular}
\end{center}

The last row is worth noting: streaming with \co{bin = 42} produces the same
averaging as the recommended Method A operating point, at the same rate, but with
100\% duty cycle instead of 21\% --- i.e. \emph{if the streaming noise excess were
resolved}, streaming would dominate Method A at every rate.

\subsection{\texorpdfstring{\co{STREAM_DURATION_SEC}, \co{total_reps} and \co{cfg.reps} --- all the same number}{STREAM\_DURATION\_SEC, total\_reps and cfg.reps}}

The tProc loops \co{cfg.reps} times and then ends. The rep count \emph{is} the
run length, it is compiled into the program, and it cannot be changed once the
program is uploaded. So the cell computes

\[
\co{total_reps} = \operatorname{round}\!\left(\frac{\co{STREAM_DURATION_SEC}}{t_\text{rep}}\right)
\]

\noindent and builds the program with \co{cfg.reps} set to it.
\co{stream()} raises \co{ValueError} if the two disagree rather than quietly
streaming a different length than you asked for. At \SI{30}{\second} and
$t_\text{rep}=\SI{240}{\micro\second}$ that is \num{125000} reps, giving
$\num{125000}/4 = \num{31250}$ samples --- which is exactly the row count of
every streamed CSV in the 2026-08-17 session.

There is no way to stop early and keep a valid file beyond interrupting the cell:
Ctrl-C ends the generator, the \co{finally} leaves the board idle, and
\co{stream_qc()["aborted"]} records that the run was short.

\subsection{\texorpdfstring{\co{poll_timeout_s} --- and why it is currently unreachable}{poll\_timeout\_s}}

\co{stream(poll_timeout_s=10.0)} is a client-side guard: if the board sends
nothing for this long, raise instead of waiting forever. It only does anything
when the poll is bounded, and as of the 2026-08-17 fix the poll is blocking by
default (\co{POLL_TIMEOUT_S = None}), so on the healthy path this guard is never
reached. That is deliberate, and \S\ref{sec:wedge} is the argument for it. What
protects an interrupted run is the generator's \co{finally}, not this timeout.

To arm it for a diagnostic run, set the class attribute for the session:

\begin{center}\co{MultipointLockinODMR.POLL_TIMEOUT_S = 5.0}\end{center}

\subsection{Stride, packet size and \texorpdfstring{\co{avg_maxlen}}{avg\_maxlen}}
\label{sec:stride}

None of these are host-side settings; they are consequences of the firmware and
the board worker, and they are what \co{stream_headroom()} reports.

\begin{center}
\small
\begin{tabular}{@{}l>{\raggedright\arraybackslash}p{10.4cm}@{}}
\toprule
\hd{Quantity} & \hd{What it is} \\
\midrule
\co{avg_maxlen} & Length of the accumulated buffer, in entries, from the
  firmware. One entry per readout per rep. \\
\co{reads_per_shot} & $n_\text{slots}\times n_\text{readouts per slot}$ --- 2 for
  the two-point default, \textbf{32} for a 16-point run with references. \\
\co{shots_that_fit} & \co{avg_maxlen // reads_per_shot}. The buffer budget is
  shared across every readout in a shot, so a 16-point run has sixteen times less
  room than a two-point one. \\
stride & \co{max(1, int(0.1 * shots_that_fit))}. The board worker transfers this
  many shots at a time --- 10\% of the buffer, chosen by qick, not by us. \\
packet size & Whatever the worker had ready when the host polled: on 2026-08-14,
  38--500 reps, median 205. It varies, and \emph{it does not matter}: the rep axis
  is carried in \co{first_rep}/\co{n_reps}, so the timestamps are exact regardless
  of how the data was parcelled up. \co{stream_qc()} proves that on the data. \\
\bottomrule
\end{tabular}
\end{center}

The failure mode is overflow: if unread samples ever reach \co{avg_maxlen}, the
worker raises and the run dies. Run \co{stream_headroom()} before a long stream
and check that \co{shots_that_fit} is several strides' worth. This has never bound
on the two-point path and is a live constraint on the multipoint one
(Ch.~\ref{ch:mpstream}).

\begin{openbox}[\co{stream_headroom()} silently failed for the whole 2026-08-17 session]
It read \co{qd.soc['readouts'][ch]['avg_maxlen']}, and \co{qd.soc} is a Pyro
\co{Proxy}: a proxy forwards \emph{method calls} but not \co{__getitem__}, so
this raised \co{'Proxy' object is not subscriptable} and the cell printed
``could not read avg_maxlen'' instead of a number. The subscriptable object is
\co{qd.soccfg}, the local \co{QickConfig} that qick's \co{make_proxy()} builds
from \co{soc.get_cfg()}. Fixed 2026-08-17; the check now runs.
\end{openbox}

\section{Verification: the guarantees, checked on the data}

Streaming's value is its exactness, so Step 5C checks it rather than assuming it.
On the 2026-08-14 run:

\begin{center}
\small
\begin{tabular}{@{}lrl@{}}
\toprule
\hd{Check} & \hd{Result} & \hd{Meaning} \\
\midrule
\co{rep_index} gaps & \textbf{0} & no reps dropped; cadence timestamps valid throughout \\
\co{rep_index} step & 4, uniformly & host binning consistent, no straddled bins \\
Duplicate rows & \textbf{0} of 31\,250 & the replay race really is absent \\
Reps collected & 125\,000 / 125\,000 & nothing lost \\
Packets & 152 & 38--500 reps each, median 205 \\
Duty cycle & 100.0\% & the FPGA never idled \\
\bottomrule
\end{tabular}
\end{center}

\co{stream()} also now refuses to reshape a packet whose payload size does not
match its own rep count. Such a packet would silently misalign every subsequent
sample and mix the two parked frequencies together --- which would look exactly
like a flat broadband noise excess. No occurrence has been observed, but it would
previously have been invisible.

\section{The buffer-overflow constraint}

The board worker transfers in strides of roughly 10\% of the accumulated buffer and
raises if unread samples reach \co{avg_maxlen}. \co{stream_headroom()} reports the
actual margin at runtime rather than assuming it:

\begin{center}
\co{\{avg_maxlen, reads_per_shot, shots_that_fit, seconds_that_fit, worker_stride_shots\}}
\end{center}

At $\tau=\SI{120}{\micro\second}$ with one read per shot the cadence is
$\sim$4\,kHz, so even a modest buffer allows a poll interval of order a second.
No overflow has been observed, and the 500-sample first packet is no noisier than
the 205-sample ones, which argues against overflow being behind the noise excess.

\section{The PL droop}
\label{sec:droop}

\begin{openbox}[Real, and specific to this mode]
With no gap between acquires the sample heats. On the 2026-08-14 run the PL fell
\textbf{25\% over 30\,s}, smoothly, from 993 to 744 ADC counts.

It is mostly common-mode and largely cancels in $z_+-z_-$: it is 85\% of the
\emph{common-mode} variance but only \textbf{2.8\%} of the $\Delta f$ variance. But
it walks the operating point down the flank, so the slopes the conversion assumes
drift out of validity over the run.

Step 5C detrends it (rolling median, corner $\sim f_s/\text{window}$) and reports
the magnitude. The proper fix is a periodic re-zero during the run, which is not
implemented.
\end{openbox}

\section{When to use it}

For anything that needs a \emph{continuous} record above a few hundred Hz --- MCG,
drone transients --- streaming is the only method that provides it: no gaps, exact
timestamps, 100\% duty cycle. Its rate is proven.

Its noise floor is currently $\sim$6$\times$ worse than burst, which for absolute
sensitivity work makes burst the better choice today. \textbf{Resolving
\S\ref{sec:streamnoise} is the single highest-value open item in this document},
because it would make streaming better than both other methods at every operating
point.
